%-------------------------------------------
% Intro page for the ERL version
%-------------------------------------------

\thispagestyle{empty}
%\begin{figure}[t]
%  \vspace*{1cm}    
%  \includegraphics[height=4.1cm]{./fig/logos/E-Robotics_league_logo.pdf}
%\end{figure}
\textbf{\center \Large 
 		Introduction}
 		
%\begin{center}
%%	\textbf{\large 
%% 		Robot Competitions Kick Innovation\\
%% 		in Cognitive Systems and Robotics\\
%% 		FP7-ICT-601012\\
%%	}
%	
%	\vspace*{10mm}
%	\hrulefill
%	\vspace*{5mm}
%
%	\textbf{%
%		{\Huge \roah}\\[1.5ex]
%		{\Large -- A Competition for Domestic Service Robots --}\\[1ex]
%	}
%
%	\hrulefill
%	\vspace*{10mm}
%\end{center}
\vspace*{0.8cm}	
\noindent 
%\hrulefill
\vspace*{0.2cm}
\\
\noindent
The European Robotic League (ERL) is an innovative concept for robot competitions. The ERL is composed of multiple ‘’Local tournaments’’, held in different research labs across Europe, with certified testbeds, and a few competitions as part of ‘’Major tournaments’’ such as RoboCup. Teams participate in a minimum of 2 tournaments (Local and/or Major) per year and get scores based on their performances. A final end of year score is computed for each team, per Task Benchmark (TBM) and Functionality Benchmark (FBM), using the best two participation in tournaments, and teams are ranked based on their final score. Prizes for the top teams (per TBM and FBM where they have participated) will be awarded during the next year’s European Robotics Forum (ERF).  For more information, please visit:

\begin{center}
\url{http://sparc-robotics.eu/the-european-robotics-league/}
\end{center}

\noindent
The \textbf{\erlir} tournament stems from its predecessor, the \textbf{RoCKIn@Work} competition, and focuses on the domain of industrial robotics in the Factory of the Future and also deals with modern automation issues. 
This document describes the rules of the \erlir competition which is in fact an updated version of the latest rules of the RoCKIn@Work\footnote{\url{http://rockinrobotchallenge.eu/rockin_d2.1.6.pdf}} 2015 competition. One task benchmark, "Fill a box with parts for manual assembly", has been modified in such a way that it can now be shared with the \textbf{RoboCup@Work} "Basic Transportation Test". An additional functionality benchmark "Navigation" has been added based on the ERL Service Robots rule book, which can be shared with RoboCup@Work "Basic Navigation Test".

\vspace{\fill}
\noindent
\small Note: This version of the ERL-IR rule book is valid until 28.02.2017.
%\hrulefill

%\begin{figure*}[b]
%	\centering
%	\includegraphics[height=\institutionLogoHeight]{./fig/logos/consortiumLogo.pdf}
%\end{figure*}
